\section{Study Design}
\label{sec:study-design}

\subsection{Research Questions}
\label{sec:rq}

We reproduce the project's frozen research questions verbatim (see
\texttt{docs/RESEARCH\_QUESTIONS.md} for the full changelog):

\begin{description}
  \item[RQ1.] How stable are scheduler rankings across independent
    workload sources?
  \item[RQ2.] How stable are scheduler rankings under temporal, provider,
    and domain shifts?
  \item[RQ3.] To what extent do rankings obtained on synthetic stress
    workloads transfer to rankings on independent real-trace-derived
    workloads?
  \item[RQ4.] Which source-native observable workload characteristics are
    associated with cross-distribution scheduler rank reversals?
  \item[RQ5.] How many independent workload windows are required before a
    comparative scheduler ranking becomes statistically stable?
  \item[RQ6.] Do representative simulated scheduler rankings and
    cross-workload rank reversals reproduce on a real serving engine?
\end{description}

Load-level dependence is a secondary robustness/sensitivity analysis applied
to RQ1--RQ4's findings, not a headline research question
(\S\ref{sec:evidence-boundaries}).

\subsection{Study Overview}
\label{sec:overview}

LSSP evaluates a fixed panel of 13 scheduling policies (\S\ref{sec:panel})
against a frozen corpus of 120 workload windows drawn from three
independently sourced traces (\S\ref{sec:sources}), across a calibrated
six-region grid of operating regions (\S\ref{sec:calibration}) and a fixed
set of primary and completion-conditioned metrics (\S\ref{sec:metrics}).
Every policy is run against every window under every operating region,
with a deterministic rep0/rep1 verification pass (\S\ref{sec:paired}),
producing $120 \times 6 \times 13 \times 2 = 18{,}720$ completed,
independently validated cells, from which per-source, per-region, and
per-metric scheduler rankings are derived. Ranking-portability statistics (Kendall's tau-b, Spearman's
rho, top-\emph{k} overlap, and pairwise reversal detection) are then
computed across pairs of sources, regions, and metrics, using workload
window as the statistical unit (\S\ref{sec:unit}). This design was frozen
(\S\ref{sec:prereg}) before any comparative scheduler outcome existed.

\subsection{Evidence Independence / Publication Boundaries}
\label{sec:evidence-boundaries}

This project shares a lineage and a simulator codebase with two other,
separately scoped manuscripts, and its design was explicitly audited for
overlap against both before this study's protocol was frozen
(\texttt{docs/OVERLAP\_LEDGER.md}, \texttt{docs/GO\_NO\_GO\_GATES.md} Gate A,
\texttt{docs/EVIDENCE\_INDEPENDENCE\_PLAN.md}):

\begin{itemize}
  \item \textbf{``The Exploitability Gap in LLM-Serving Scheduler
    Portfolios'' (LLM 2026).} That manuscript's
    \texttt{public\_replay\_load\_scaling\_v1/v2} experiment used a fixed set
    of 60 canonical public-trace windows (20 BurstGPT + 20 Azure-2023
    conversation + 20 Azure-2023 code), a load-factor grid
    $\{1,2,4,8,16,32,64,128\}$, an 8-policy portfolio, and ANWG/SBS-VBS
    exploitability diagnostics. That result is treated here as
    \texttt{PRIOR\_RESULT\_REFERENCE\_ONLY}: citable as prior motivation, never
    restated as this project's own evidence. This project's own window
    sampling is independently constructed (\S\ref{sec:windows}) and
    verified non-identical to that 60-window set; load-level dependence is
    accordingly demoted to a secondary analysis (\S\ref{sec:rq}) using this
    project's own calibration protocol (\S\ref{sec:calibration}), never that
    load grid.
  \item \textbf{``Module Counterfactuals for Attributing LLM Serving
    Scheduler Effects'' (a separate module-intervention manuscript).} This
    study makes no module-counterfactual or module-attribution claim; RQ4's
    workload-descriptor analysis (\S\ref{sec:methodology}) is an offline,
    descriptive association analysis, not a deployed selector or an
    intervention-based attribution method.
  \item Other project lineage (an unrelated database/systems research
    thread) is unconnected to this manuscript's content and is not
    referenced here.
\end{itemize}

\subsection{Experimental Unit and Paired Design}
\label{sec:unit}

The statistical unit of analysis is the workload window: a fixed-size
(200-request) contiguous slice of a source trace, uniquely identified and
frozen before any scheduler outcome existed (\S\ref{sec:windows}). Every
policy in the panel (\S\ref{sec:panel}) is run against every window under
every calibrated operating region (\S\ref{sec:calibration}), so that
ranking-portability statistics are computed on genuinely paired
observations -- the same window, region, and metric definition, varying
only the policy -- rather than on separately tuned per-source experiments.
Source, operating region, and metric are treated as factors of interest;
policy identity is the object whose relative ordering is measured.

\subsection{Preregistration}
\label{sec:prereg}

The workload window set (120 windows: 40 per source, 10 Stage-0-reused +
30 Pilot-V2-new per source; \S\ref{sec:windows}), the scheduler panel and
its selection rule (\S\ref{sec:panel}), the metric contract
(\S\ref{sec:metrics}), the six-region operating-load calibration
(\S\ref{sec:calibration}), and the ranking-portability analysis plan
(reversal-detection thresholds, bootstrap procedure, sample-complexity
design) were each frozen, and cross-checked against a structural validation
gate, before any Pilot-V2 comparative scheduler outcome was generated.
\texttt{docs/RANKING\_PORTABILITY\_PILOT\_V2\_PROTOCOL.md},
\texttt{docs/RANKING\_PORTABILITY\_ANALYSIS\_PLAN.md}, and
\texttt{docs/RANKING\_PORTABILITY\_WINDOW\_FREEZE.md} record the exact frozen
content and the gate checks each passed.
